\chapter{Possible Extensions}
\label{chap:disc}



\paragraph{Extension 1} Since the short rate is determined by the latent state, $r(t)=\tilde{r}(z_t)$, its precision depends on how well the encoder identifies $z_t$. Incorporating inflation information is therefore not about directly refining the short rate, but about improving state identification. This motivates extending the model to include inflation swap data as additional input to the encoder. (objective is here not to fit better, but to predict better)

Makroøkonomisk data man kan overveje at inddrage i input matricen:
\begin{itemize}
    \item Realized volatility of all swap rates
    \item Inflation swap rates
    \item Stock prices in the different currencies
    \item volatility of stock prices
\end{itemize}




\paragraph{Extension 2} While PCA provides a useful low-dimensional statistical benchmark, it does not account for dynamic state estimation under a structural term-structure model. 
The Kalman filter, by contrast, is the standard econometric method for estimating latent factors in affine models, combining pricing restrictions with time-series dynamics. Comparing Rolf's model to a Kalman filter benchmark therefore allows us to assess whether the neutral architecture improves state identification and pricing perfomance relative to a well-established filtering approach.



\paragraph{Argument for Extension 2} As a first step, we establish a Kalman filter benchmark, since it represents the standard state-space approach to estimating latent term-structure factors under measurement noise and dynamic evolution. By comparing the encoder-implied states in Rolf’s model with Kalman-filtered factors, we can assess differences in stability, robustness, and out-of-sample performance. This benchmark will guide whether further improvements should focus on filtering, hybridization, or volatility specification within the neural framework.








Noter
\begin{itemize}
    \item usikkerhed driver større risikopræmier. Dette har vi set i det seneste pr med den Japanske yenn, vedrørende alt den usikkerhed der har være i verdenen med både krig samt Trumps beslutninger truffet på USAs vegne.
    \item Er modellen stiafhængig? Nej, det er den ikke i den klassiske forstand.
    \item \textbf{Kalman filter vs. Rolfsmetode} Fordele og ulemper man behøver ikke i vores at antage en eller anden sammenhæng, til gengæld kan man ikke nedlægge en økonomisk teori ned over.
\end{itemize}


\section{Implementation of the two extensions}


\subsection{Extension 1}


else try to extend our model to also account for the swap rates realized volatility in order to see if it makes our reconstruction error better. write this mathematically how we can extend this in our model.

Realiseret volatiltet på swaprenterne, og indsætte dette som node, som det første

for hver swaprate og hver tenor fås et punkt






\subsection{Extension 2}

The Kalman filter takes noisy observed data and produces the best (Gaussian-optimal) estimate of the hidden latent state that evolves over time.

We observe swap rates $S_t$, and we assume that they are driven by a small number of latent factors $x_t$. Those factors evolve dynamically

$$x_{t+1}=A x_t+\eta_t$$

the observed swaps are noisy measurements of those factors
$$S_t=Hx_t+\varepsilon_t$$

The Kalman filter then 
\begin{enumerate}
    \item \textbf{Predicts} today' state using yesterday's state.
    \item \textbf{Updates} that prediction using today's observed swap curve.
    \item Repeats recursively through time
\end{enumerate}

In contrast, Rolf's encoder estimates $z_t$ from each day independently, witouht filtering over time.

So the question when comparing these two are:

Does Rolf’s encoder produce latent states that are as informative and stable as Kalman-filtered states when both are inferred from the same observed swap curve?






\subsection{Data}


Data is the quotes on interest rate options, specifically a swaption, on the time period 1990 to february 2026. For the currencys USD, EUR and JPY.







Euro: 2010 - end date (EUSA) OIS swaps, broker ICAP, ICPL
Old Euro
USD: 


Arguement for volatility


MERGE:
US SWAPVOL OIS 2023-2025
US SWAPVOL LIB 2001-2023 

MERGE:
EURO SWAPVOL OIS 2012-2026
EURO SWAPVOL LIB fra 2001-2012



USD
EUR
JPY
DKK
NOK
SEK



